\section{PMI and PMI-B budgets}

This chapter develops the PMI identity and the PMI-B budget layer in
full.  The layer converts the exact word-orbit identity into counting
inequalities that constrain the step weights of a closed cycle word.

\subsection{The margin}

For a word $w=(t_0,\dots,t_{L-1})$ with prefix weights
$\wordWeight_j$, define the margin after $j$ steps by
\begin{equation*}
  m_j=j\log_2 5-\wordWeight_j.
\end{equation*}
The margin measures how much 2-adic budget remains after $j$ exact
steps.  A prefix with $m_j\ge0$ is called a \emph{nonnegative-margin
prefix}; a proper prefix with $m_j<0$ is called a \emph{bad prefix}.

\subsection{The PMI identity}

\begin{theorem}[PMI identity]
For every valid word $w$ of length $L$,
\begin{equation*}
  \sum_{j=0}^{L-1}2^{-m_j}
    =\frac{5\,\wordA(w)}{5^L}.
\end{equation*}
\end{theorem}

\begin{proof}
Write the word molecule as
\begin{equation*}
  \wordA(w)=\sum_{j=0}^{L-1}
    2^{\wordWeight_j}\,5^{L-1-j}.
\end{equation*}
Then
\begin{equation*}
  \frac{5\,\wordA(w)}{5^L}
    =\sum_{j=0}^{L-1}2^{\wordWeight_j}\,5^{-j}.
\end{equation*}
Since
\begin{equation*}
  2^{-m_j}=2^{\wordWeight_j-j\log_2 5}
    =2^{\wordWeight_j}\,5^{-j},
\end{equation*}
the identity follows.  The cleared-denominator form is
\begin{equation*}
  \sum_{j=0}^{L-1}2^{\wordWeight_j}\,5^{L-j}
    =5\,\wordA(w).
\end{equation*}
\end{proof}

\subsection{The cycle form}

Let $w$ be a closed cycle word of period $P$ with cycle state $m$,
total weight $T=\wordWeight(w)$, and prefix weights
$\wordWeight_j$.  Define
\begin{equation*}
  \mathrm{aTotal5}(P,t)
    =\sum_{j=0}^{P-1}2^{\wordWeight_j}\,5^{P-j}.
\end{equation*}
The closedness
\begin{equation*}
  \wordOrbit(w,m)=m
\end{equation*}
gives the cycle equation
\begin{equation*}
  2^T\,m=5^P\,m+\wordA(w).
\end{equation*}
Substituting into the cleared form yields
\begin{equation*}
  \mathrm{aTotal5}(P,t)
    =5m\left(2^T-5^P\right).
\end{equation*}
This is the algebraic form
\texttt{StringFlow.PMI.pmi\_\allowbreak algebraic} in Lean.

\subsection{The counting bound}

Write $B$ for the number of bad proper prefixes of a closed cycle
word of period $P$.  The $j=0$ term of $\mathrm{aTotal5}$ contributes
\begin{equation*}
  2^{\wordWeight_0}\,5^{P}=5^P.
\end{equation*}
Every bad prefix $j$ satisfies
\begin{equation*}
  \wordWeight_j>j\log_2 5,
\end{equation*}
equivalently
\begin{equation*}
  2^{\wordWeight_j}>5^j,
\end{equation*}
so its contribution satisfies
\begin{equation*}
  2^{\wordWeight_j}\,5^{P-j}>5^P.
\end{equation*}
Combining the $j=0$ term with the bad-prefix contributions gives
\begin{equation*}
  5^P+B\cdot5^P<\mathrm{aTotal5}(P,t).
\end{equation*}
The paper writes this with $\le$; the strict form follows because
bad prefixes are defined by strict margin violation.  Using the cycle
form,
\begin{equation*}
  5^P+B\cdot5^P<5m\left(2^T-5^P\right).
\end{equation*}

\subsection{Small budget and no bad prefix}

\begin{corollary}
If
\begin{equation*}
  5m\left(2^T-5^P\right)<2\cdot5^P,
\end{equation*}
then $B=0$: every proper prefix satisfies
\begin{equation*}
  2^{\wordWeight_j}<5^j.
\end{equation*}
\end{corollary}

\begin{proof}
The counting bound gives
\begin{equation*}
  B\cdot5^P<5m\left(2^T-5^P\right)-5^P<5^P,
\end{equation*}
so $B<1$, hence $B=0$.
\end{proof}

\subsection{Margin balance}

Let $C_1,C_2,C_3$ be the numbers of steps of weight $1$, weight $2$,
and weight at least $3$ in a word, and put $\alpha=\log_2 5$.  The
nonnegative-margin prefix condition implies
\begin{equation*}
  (3-\alpha)C_3\le(\alpha-1)C_1+(\alpha-2)C_2.
\end{equation*}
The strict version, used for closed cycle words, replaces $\le$ by
$<$.

The balance records that a C3 step spends at least $3-\alpha$ units of
margin, while a weight-$1$ step earns $\alpha-1$ and a weight-$2$
step earns $\alpha-2$.  The proof telescopes the margin differences
along the word and compares the final margin with the total weight.
The Lean statements are
\texttt{cycleWord\_\allowbreak margin\_\allowbreak balance} and
\texttt{cycleWord\_\allowbreak margin\_\allowbreak balance\_\allowbreak
strict}.

\subsection{Worked example: the exact prefix word}

For the exact prefix word
\begin{equation*}
  w=[2,1,2]
\end{equation*}
from $7$, we have $L=3$, total weight $T=5$, and
$\wordA(w)=53$, with
\begin{equation*}
  2^5\cdot29=928=5^3\cdot7+53.
\end{equation*}
The prefix weights are
\begin{equation*}
  \wordWeight_0=0,\quad \wordWeight_1=2,\quad \wordWeight_2=3,
\end{equation*}
so the margins are
\begin{equation*}
  m_0=0,\qquad
  m_1=\log_2 5-2\approx0.3219,\qquad
  m_2=2\log_2 5-3\approx1.6439.
\end{equation*}
All proper prefixes are nonnegative, and the PMI sum is
\begin{equation*}
  \sum_{j=0}^{2}2^{-m_j}
    =1+\frac45+\frac8{25}
    =\frac{53}{25}
    =\frac{5\wordA(w)}{5^3}.
\end{equation*}

\subsection{Worked example: a bad prefix}

The word $[3,1]$ has prefix weights
\begin{equation*}
  \wordWeight_0=0,\quad \wordWeight_1=3,\quad \wordWeight_2=4,
\end{equation*}
and molecule
\begin{equation*}
  \wordA([3,1])=2^0\cdot5^1+2^3\cdot5^0=13.
\end{equation*}
The margins are
\begin{equation*}
  m_0=0,\qquad
  m_1=\log_2 5-3\approx-0.6781,\qquad
  m_2=2\log_2 5-4\approx0.6439.
\end{equation*}
The proper prefix of length $1$ is bad, so $B=1$.  The cleared sum is
\begin{equation*}
  \mathrm{aTotal5}(2,t)
    =2^0\cdot5^2+2^3\cdot5^1
    =25+40=65=5\wordA([3,1]).
\end{equation*}
The counting bound
\begin{equation*}
  5^2+B\cdot5^2=25+25=50<65
\end{equation*}
holds strictly.  This example is the sharpest elementary illustration
of the PMI-B counting mechanism: one bad prefix consumes one extra
unit of $5^P$, and the algebra keeps the exact count.

\subsection{Why the budget is exact}

No probabilistic replacement is involved.  The margin is defined by
the exact prefix weights of the word, and every inequality in this
chapter is derived from the identity
\begin{equation*}
  2^{\wordWeight_j}\,5^{L-j}=5^{L-m_j}.
\end{equation*}
The constants $\alpha$, $3-\alpha$, and the threshold $2\cdot5^P$
appear because the word algebra is exact, not because of a statistical
model.

\subsection{Lean statements}

\begin{itemize}[leftmargin=*]
\item \texttt{StringFlow.PMI.aTotal5}: the cleared sum.
\item the identity $\mathrm{aTotal5}=5\,\mathrm{aTotal}$ in
  \texttt{Pmi.lean}: the algebraic clearing.
\item \texttt{StringFlow.PMI.pmi\_\allowbreak algebraic}: the cycle form
  $\mathrm{aTotal5}(P,t)=5m(2^T-5^P)$.
\item \texttt{pmi\_\allowbreak b\_\allowbreak count}: the counting
  bound.
\item \texttt{pmi\_\allowbreak b\_\allowbreak no\_\allowbreak
  bad\_\allowbreak prefix}: the small-budget corollary.
\item \texttt{cycleWord\_\allowbreak margin\_\allowbreak balance} and
  \texttt{cycleWord\_\allowbreak margin\_\allowbreak balance\_\allowbreak
  strict}: the margin balance.
\end{itemize}

\subsection{Status}

\begin{center}
\small
\begin{tabular}{@{}ll@{}}
\toprule
Statement & Status \\
\midrule
PMI identity & Lean: compiled \\
cycle form & Lean: compiled \\
PMI-B counting bound & Lean: compiled \\
small-budget corollary & Lean: compiled \\
margin balance & Lean: compiled \\
\bottomrule
\end{tabular}
\end{center}
